\section{Local inverses and walk expansions}\label{sec:local-walks}

\subsection{Averaging penalties and the inverse at zero field}
\label{subsec:blocked-zero-inverse}

Fix an odd $L\ge5$ and put $a=\eta_k=L^{-k}$. Write $e_k$ for the
charge $e_a$ at this mesh. The level data $k,\mu,R,e_k$ are independent,
with $0<e_k\le e_0$.
Unit blocks have side
one, and the next blocks have side $L$. In this and the following
subsections a finite period $R$ is a positive integer multiple of $L$.
For $z\in\Z^3$ the unit block is
$\Delta_z=\{c\in a\Z^3:|c-z|_\infty<1/2\}$.
For $z'\in L\Z^3$ the $L$-block is
$E_{z'}=\{c\in a\Z^3:|c-z'|_\infty<L/2\}$.
These partitions are taken modulo $R\Z^3$ on the torus. Since $L$
is odd, no fine site lies on a block boundary. Their means, including
the cluster mean, are denoted $\Pi_1$ and $\Pi_L$. At nonzero connection
the means use the gauge phases of \cref{sec:introduction}.
With the constants and coefficients defined there, the blocked operators
at zero connection are
\begin{equation}\label{eq:blocked-zero-operator}
 F_y(0)=D_0+\mu+\alpha_y\Pi_1+\beta_y\Pi_L\quad(k\ge1),
 \qquad
 F_y(0)=D_0+\mu+b_L\Pi_L+i\gamma_5y\quad(k=0).
\end{equation}
Here $0\le\mu\le1$. In the first formula $y\in\R\cup\{\infty\}$;
in the second $y\in\R$. We use $-\infty=\infty$ when an adjoint
identity reverses this parameter. The notation $F_y$ denotes an operator in this
section, not the curvature.

The penalty bounds follow by diagonalizing $\gamma_5$. On its eigenspace
of sign $\epsilon$, the two eigenvalues and their sum are
\[
 \frac{b_ki\epsilon y}{b_k+i\epsilon y},\quad
 \frac{b_k^2b_L}{(b_k+b_L+i\epsilon y)(b_k+i\epsilon y)},\quad
 \frac{b_k(b_L+i\epsilon y)}{b_k+b_L+i\epsilon y}.
\]
For the first eigenvalue $v_1$ and the sum $v_2$, direct calculation gives
\[
 |v_1|^2=b_k\operatorname{Re}v_1,\qquad
 |v_2|^2\le b_k\operatorname{Re}v_2,\qquad
 \operatorname{Re}v_2\ge\frac{b_kb_L}{b_k+b_L}\ge\frac{b_-}{L}.
\]
Also $|\alpha_y|\le b_k\le b$ and $|\beta_y|\le b_L$.
On an $L$-block the decomposition
$\operatorname{Pen}=\alpha_y(\Pi_1-\Pi_L)+(\alpha_y+\beta_y)\Pi_L$
has orthogonal ranges. The scalar inequalities on the two spin
eigenspaces therefore give
\begin{equation}\label{eq:penalty-estimates}
 \operatorname{Re}(f,\operatorname{Pen}f)_E
      \ge\frac{b_-}{L}\|P_Ef\|_{2,E}^2,
 \qquad
 \|\operatorname{Pen}f\|_{2,E}^2
      \le b\operatorname{Re}(f,\operatorname{Pen}f)_E.
\end{equation}
The kernel on $E$ is
$\alpha_y\mathbf1[c,c'\text{ in the same unit block}]+\beta_yL^{-3}$,
and its norm is at most $2b$.

\begin{proposition}\label{prop:zero-blocked-inverse}
There are $C_I,C_5,\gamma_5'>0$, depending on $L,\rho$ and a fixed
$\alpha\in(0,1)$ only, such that $F_y(0)$ is invertible and
\[
 \|F_y(0)^{-1}\|_{2\to2}\le LC_I\qquad(k\ge1).
\]
If $f$ is supported over one unit block $\Delta_{z'}$ and
$v=F_y(0)^{-1}f$, then for $x$ over $\Delta_z$ and
$0<d_X(x,\widetilde x)\le1$,
\begin{equation}\label{eq:zero-blocked-holder}
 |v(x)|_2+d_X(x,\widetilde x)^{-\alpha}|v(x)-v(\widetilde x)|_2
 \le C_5\exp(-\gamma_5'd(z,z'))\|f\|_\infty.
\end{equation}
This includes pairs in one cluster. The constants are uniform in
$k,R$, the character, $\mu$, and $y$.
\end{proposition}

\begin{proof}
First suppose $k\ge1$. Put $f=F_y(0)u$ and
$X=\operatorname{Re}(u,f)$. The real-part identity
and \eqref{eq:penalty-estimates} imply
\[
 \sum_E\|P_Eu\|_{2,E}^2\le\frac L{b_-}X,\qquad
 \|\operatorname{Pen}u\|_2^2\le bX,\qquad
 \|\mu u\|_2^2\le X.
\]
Therefore $\|D_0u\|_2\le\|f\|_2+(1+\sqrt b)X^{1/2}$.
Use \eqref{eq:free-coercivity} for $D_0u$ and
\eqref{eq:chamber-poincare} on each $L$-block. With
$K_1=\sqrt L+\sqrt{Lb}$ this gives
\[
 \|u\|_2^2\le L\left(\frac1{b_-}+\frac{2C_PK_1^2}{c_b}\right)
           \|u\|_2\|f\|_2
       +\frac{2C_P}{c_b}L^2\|f\|_2^2.
\]
Hence $\|u\|_2\le LC_I\|f\|_2$, where
$C_I=b_-^{-1}+2C_PK_1^2/c_b+(2C_P/c_b)^{1/2}$.
The identity $F_y(0)^*=\gamma_5F_{-y}(0)\gamma_5$ gives the same
lower bound for the adjoint. Thus $F_y(0)$ has closed, dense range
and is onto, including on the infinite lattice.

To obtain decay, let $\psi$ be a bounded real periodic function with
$|\psi(c)-\psi(c')|\le a_0|c-c'|_1/L$. The mass commutes with this
weight. The external hopping term gives
$\|\exp(\psi)[D_0,\exp(-\psi)]\|_{2\to2}
\le3(1+\rho)\exp(1)a_0/L$.
The penalty kernel is supported in one $L$-block, has norm at most
$2b$, and $|\psi(c)-\psi(c')|\le3a_0$ there. Its contribution is at
most $6L^4b\exp(1)a_0/L$. Put
\[
 c_w=\exp(1)(3(1+\rho)+6L^4b),\qquad
 a_0=\min(1/3,(2C_Ic_w)^{-1}).
\]
Applying the $L^2$ inverse bound and absorbing the commutator gives
\begin{equation}\label{eq:weighted-block-inverse}
 \|\exp(\psi)F_y(0)^{-1}f\|_2
       \le2LC_I\|\exp(\psi)f\|_2.
\end{equation}
Take truncated multiples of the $\ell^1$ distance, divided by $L$,
to the source $L$-block. Letting the truncation increase proves an
exponentially weighted $L^2$ bound.

The passage from $L^2$ to the sup norm uses the free kernel. On the
covering lattice choose $c^\circ\in a\Z^3$ and $r\ge8a$, and put
\[
 B=\{|c-c^\circ|_\infty<r\},\quad
 B^+=\{|c-c^\circ|_\infty<r+a\},\quad
 B'=\{|c-c^\circ|_\infty\le r/4\}.
\]
Suppose $u$ is defined over $B^+$ and $(D_0+\mu)u=h$ over $B$,
with $0\le\mu a\le1$. Both evaluation points below lie over $B'$,
and $C$ depends on the fixed $\alpha\in(0,1)$ and $\rho$.
Let $\chi$ equal one
on the concentric half cube and zero outside the cube, with Lipschitz
constant $2/r$. Then
\[
 \chi u=G(\chi h)+G([D_0,\chi]u).
\]
The commutator contains only the external hop. Its rows are supported
where $|c-c^\circ|_\infty>r/2-a$. It satisfies
$\|[D_0,\chi]u\|_1\le45r^{1/2}\|u\|_{2,B^+}$.
For a point of the quarter cube, this support is at distance at least
$r/4-a\ge r/8$. The bounds \eqref{eq:mesh-free-kernel} and
\eqref{eq:free-holder-kernel} consequently give
\begin{equation}\label{eq:interior-dirac-estimate}
 \begin{split}
 |u(x)|_2&\le C\big(r\|h\|_{\infty,B}+r^{-3/2}\|u\|_{2,B^+}\big),\\
 |u(x)-u(\widetilde x)|_2&\le
 C\left(\frac{d_X(x,\widetilde x)}r\right)^\alpha
       \big(r\|h\|_{\infty,B}+r^{-3/2}\|u\|_{2,B^+}\big).
 \end{split}
\end{equation}
For the first term, sum $(a+|c-c'|_\infty)^{-2}$ over the cube;
the result is $O(r)$. For the difference, the analogous sum with
exponent $2+\alpha$ is $O(r^{1-\alpha}/(1-\alpha))$.
The cluster average in $\dd x'$ contributes one to both sums.
For the commutator term, its support separation gives the size bound;
summing bond differences along the chamber path gives the difference
bound. These estimates explain why the constants do not grow when
$a\downarrow0$.

For $u=F_y(0)^{-1}f$, use
$h=f-\operatorname{Pen}u$ in \eqref{eq:interior-dirac-estimate}, with
$r=L/4$. The enlarged cube meets at most eight $L$-blocks. On each,
the penalty kernel and Cauchy--Schwarz give
$\|\operatorname{Pen}u\|_\infty\le2bL^{3/2}\|u\|_{2,E}$.
The weighted bound \eqref{eq:weighted-block-inverse} therefore bounds
both terms in \eqref{eq:interior-dirac-estimate} by
$CL\exp(-a_0d_1(x,E'))\|f\|_\infty$, where
$d_1(x,E')$ is the $\ell^1$ distance to the source $L$-block divided
by $L$. This applies for $k\ge2$.
For $k=1$, the pointwise $L^2$ bound costs only
$\sqrt{48}L^{3/2}$; $L$ is fixed. Larger Hölder separations are
controlled by the two sup bounds. Finally,
$d_1(x,E')\ge d(z,z')/L-3$, and the source unit block lies in $E'$.
This proves \eqref{eq:zero-blocked-holder} with
$\gamma_5'=a_0/L$. The argument uses only periodic moduli and so
applies to every unitary character.
\end{proof}

The first level has a different resolvent parameter. On an $L$-block
the real-part bound and Poincar\'e inequality imply
\[
 \operatorname{Re}(u,F_y(0)u)\ge\frac{c_r}{L}\|u\|_2^2,
 \qquad c_r=\min(b_-,\rho/(2C_PL)).
\]
The term $i\gamma_5y$ has zero real part. Conjugation by a weight with
Lipschitz constant $a_z$ changes the hopping and penalty terms by at
most $\exp(1)(6+3b)a_z$. Choose
\[
 a_z\le\min\bigl((3L)^{-1},c_r/(2\exp(1)L(6+3b))\bigr).
\]
The weighted inverse then has $L^2$ norm at most $2L/c_r$.
On the other hand its part other than $i\gamma_5y$ has norm at most
$60+b$. For $|y|\ge2(60+b)$, the inverse norm is at most $2/|y|$.
Combining the two ranges and converting to a kernel, with its factor
$48$, gives
\begin{equation}\label{eq:level-zero-kernel}
 |F_y(0)^{-1}(x,x')|_2\le C_{00}(1+|y|)^{-1}
                         \exp(-a_zd(c,c')).
\end{equation}
This factor in $y$ is needed when the determinant is represented by an
integral of resolvents.

\subsection{Small connections and local inverses}\label{subsec:local-inverses}

Use the covariant block means, short quotients, and row norms of
\cref{sec:introduction}. Replacing $D_0$ by $D_A$ and $\Pi_1,\Pi_L$
by their covariant versions in \eqref{eq:blocked-zero-operator}
gives $F_y(A)$. For real
$A$ these idempotents are orthogonal projections.
Inserting the block partition between two factors and
using the triangle inequality for $d$ proves
$N_\gamma(TT')\le N_\gamma(T)N_\gamma(T')$.
Enlarge $C_5,C_{00}$ to at least one, and choose
\[
 \gamma=\tfrac12\min(\gamma_5',a_z),\qquad
 a'=\max(C_5,C_{00})\sum_{n\in\Z^3}\exp(-\gamma|n|_\infty).
\]
The preceding estimates imply
\begin{equation}\label{eq:zero-row-bounds}
 N_\gamma(F_y(0)^{-1})\le a'w_k(y),\qquad
 N_\gamma(\delta_{\alpha,0}F_y(0)^{-1})\le a'.
\end{equation}
At $k=0$ the first bound holds also in $N_\gamma^{\mathrm{ker}}$.

\begin{proposition}\label{prop:small-connection-inverse}
Let $B$ be a complex connection with $e_k\|B\|_\infty\le\delta$, where
$3L\delta\le1$ and $2C_Va'\delta\le1$, with
$C_V=2\exp(1)(42+6b)$. Then $F_y(B)$ is invertible and
\begin{equation}\label{eq:small-row-bounds}
 N_\gamma(F_y(B)^{-1})\le2a'w_k(y),\qquad
 N_\gamma(\delta_{\alpha,B}F_y(B)^{-1})\le18a'.
\end{equation}
At $k=0$ the first bound holds also in $N_\gamma^{\mathrm{ker}}$.
\end{proposition}

\begin{proof}
Write $V=F_y(B)-F_y(0)$. The internal hops, the mass, and the
resolvent term at level zero cancel. For the external hop,
$|u_B-1|\le\eta_k\delta\exp(\eta_k\delta)$; its weighted row sum
is at most $42\exp(\gamma)\exp(\delta)\delta$. The two averaging
penalties contribute respectively
$3b\delta\exp(3\delta)$ and
$3b\delta\exp(\gamma(L-1))\exp(3L\delta)$. Indeed, the nested paths
have total length at most $3/2$, and the path to an $L$-block centre
has length at most $3L/2$. Summing their normalized block kernels
removes the numbers of sites. The bounds $\gamma L\le\log2$ and
$3L\delta\le1$ give $N_\gamma(V)\le C_V\delta$.

Thus $N_\gamma(VF_y(0)^{-1})\le1/2$. The Neumann series gives
\[
 F_y(B)^{-1}=F_y(0)^{-1}
                  \sum_{n\ge0}(-VF_y(0)^{-1})^n
             =\sum_{n\ge0}(-F_y(0)^{-1}V)^nF_y(0)^{-1}.
\]
It proves the first bound in \eqref{eq:small-row-bounds}. The first
form gives $N_\gamma(\delta_{\alpha,0}F_y(B)^{-1})\le2a'$.
For distinct photon sites in a short pair,
$d_X^{-\alpha}|U_B-1|\le3\delta\exp(3\delta)\le1/4$;
for intracluster pairs the expression is zero. Evaluation at the
second point of a short pair uses at most $27$ neighbouring blocks,
and costs at most $27\exp(\gamma)$. Its contribution to the covariant
quotient is at most $16a'$, proving the second bound. At level zero,
the second form of the Neumann series keeps the kernel norm in the
last factor, so the kernel estimate follows as well.
\end{proof}

Return to the side $\ell=R$ and the domain
$\widetilde{\mathcal G}_k$ of \cref{sec:introduction}.
On the concentric $3M$ enlargement $W_\square$ of an $M$-cube,
\begin{equation}\label{eq:local-small-gauge}
 A=B+d\lambda,\qquad e_k\|B\|_\infty<2e_k^\varepsilon,
 \qquad e_k\|\operatorname{Im}A\|_\infty<e_k^\varepsilon.
\end{equation}
This follows by restricting $B_0+A_1$ and $\lambda$ to that cube.
The domain is open: a sufficiently small complex perturbation can be
added to $A_1$. It is invariant under real gauge transformations.

For a lift $z_\square$ of the cube centre put
$W_\square=\{|c-z_\square|_\infty<3M/2\}$ and
$W_\square^-=\{|c-z_\square|_\infty<3M/2-\eta_k\}$.
Embed $W_\square$ in the positive chamber of the periodic torus of
side $6M$, whose lattice is translated by $\eta_k(1,1,1)/2$.
Its level-$j$ block centres are
$L^j\eta_k(1,1,1)/2+L^j\eta_k\Z^3$ for $0\le j\le k+1$.
Translation by $L(1,1,1)/2$ identifies these blocks with the
unshifted nested blocks, because $L^{k+1-j}$ is odd. The zero-field
and small-connection estimates therefore apply on this torus with
the same constants.
Fold its photon coordinates across the planes at multiples of $3M$.
Adjacent sites fold either to adjacent sites in $W_\square$ or to
one site; in the latter case assign phase one. Pull back $A$ by this
fold to obtain $A^W$. The fold preserves the supremum norm and
commutes with $d$ on scalar functions. It also preserves the nested
blocks, since $3M/(L^j\eta_k)$ is odd for $j\le k+1$.

Let $\mathrm{ext}$ extend a function over $W_\square$ by zero in the
seven other chambers, and let $\mathrm{res}$ restrict to its embedded
copy. Labels are unchanged in these two maps. Define
\[
 S_\square(A)=\mathrm{res}\,F_y(A^W)^{-1}\mathrm{ext}.
\]
For $A=B+d\lambda$ as in \eqref{eq:local-small-gauge},
$A^W=B^W+d\lambda^W$, with $\lambda^W$ real. Gauge conjugation by
$\exp(ie_k\lambda^W)$ reduces the inverse to the small connection
$B^W$. If $6Le_0^\varepsilon\le1$ and
$4C_Va'e_0^\varepsilon\le1$, \cref{prop:small-connection-inverse}
applies to every $0<e_k\le e_0$. Restriction, extension, and unitary
gauge factors do not increase the localized norms. Thus
\begin{equation}\label{eq:local-inverse-bounds}
 N_\gamma(S_\square(A))\le2a'w_k(y),\qquad
 N_\gamma(\delta_{\alpha,A}S_\square(A))\le18a'.
\end{equation}
The second operator here is formed on the local cube before extension
to the original torus. It does not differentiate a discontinuous
extension by zero.

The local inverse has both identities needed for patching. If $f$ is
defined over $W_\square$, then
\begin{equation}\label{eq:local-right-identity}
 F_y(A)S_\square(A)f=f\quad\text{over }W_\square^-.
\end{equation}
If $f$ is supported over $W_\square^-$, then $F_y(A)f$ is supported
over $W_\square$ and
\begin{equation}\label{eq:local-left-identity}
 S_\square(A)F_y(A)f=f.
\end{equation}
For \eqref{eq:local-right-identity}, the rows at an interior site use
only its internal neighbours, its external neighbour, and its
$L$-block. These lie in $W_\square$, where the original and folded
operators agree. For \eqref{eq:local-left-identity}, the support
margin ensures that applying $F_y$ cannot leave $W_\square$.
After embedding, applying the folded operator to the extension of
$f$ agrees with extending $F_y(A)f$. Apply the inverse and restrict.
For an antiperiodic spinor on the original torus, use one lift and
then its signed translates; the enlarged cubes do not overlap their
translates. These arguments also show that $S_\square$ reads $A$
only in $W_\square$ and is analytic in those connection variables.

\subsection{The parametrix and the walk expansion}\label{subsec:propagator-walks}

Choose the product cutoff $h_\square$ with centre $z_\square$ and
one-dimensional factor
\[
 h_0(t)=\begin{cases}
 1,&|t|\le1/4,\\
 \cos(\pi(2|t|-1/2)/2),&1/4\le|t|\le3/4,\\
 0,&|t|\ge3/4.
 \end{cases}
\]
On a lift of the cube set
\[
 h_\square(c)=\prod_{\mu=1}^3
 h_0\bigl((c_\mu-z_{\square,\mu})/M\bigr),\qquad
 h_\square(c,g)=h_\square(c).
\]
Extend this with period $R$ and set it to zero away from the periodic
images of the lifted cube neighbourhood.
The cutoffs satisfy $\sum_\square h_\square^2=1$,
$\supp h_\square\subset\{|c-z_\square|_\infty<3M/4\}$, and
$|h_\square(c)-h_\square(c')|\le\pi|c-c'|_1/M$.
Put $C_\square=-[F_y(A),h_\square]$. Internal hops commute with
$h_\square$, and the penalty has range at most $L$. The phase bounds
and the same normalized sums used in the perturbation estimate give
\begin{equation}\label{eq:cutoff-commutator-bound}
 N_\gamma(C_\square)\le\pi C_V/M.
\end{equation}
The kernel is
$(h_\square(c)-h_\square(c'))F_y(A)(x,x')$. Its two photon sites lie
within $3M/4+L$ of $z_\square$. If $M>4L$, both lie in
$W_\square^-$; multiplication by $h_{\square'}$ can give a nonzero
operator only when $d(z_\square,z_{\square'})\le M$.

Using \eqref{eq:local-right-identity} and
\eqref{eq:local-left-identity} gives
\begin{align}
 F_yh_\square S_\square h_\square
       &=h_\square^2-C_\square S_\square h_\square,
                                   \label{eq:parametrix-right}\\
 h_\square S_\square h_\square F_y
       &=h_\square^2+h_\square S_\square C_\square.
                                   \label{eq:parametrix-left}
\end{align}
Let $S^{\mathrm p}=\sum_\square h_\square S_\square h_\square$ and
$R^{\mathrm p}=\sum_\square C_\square S_\square h_\square$.
At most $27$ cubes contribute to one row of $R^{\mathrm p}$.
Consequently,
\[
 N_\gamma(R^{\mathrm p})\le q_0:=\frac{54\pi C_Va'}M.
\]
Choose $M\ge108\pi C_Va'$. Equation~\eqref{eq:parametrix-right}
now gives a right inverse by a convergent series:
\begin{equation}\label{eq:global-parametrix-inverse}
 S_y(A)=F_y(A)^{-1}=S^{\mathrm p}\sum_{n\ge0}(R^{\mathrm p})^n.
\end{equation}
The finite-dimensional operator $F_y(A)$ is therefore invertible.
The left identity gives the second form with
$R^{\mathrm p\prime}=-\sum_\square h_\square S_\square C_\square$.

The terms in \eqref{eq:global-parametrix-inverse} have a local
description. A walk of length $n$ is a sequence
$\omega=(\square_0,\ldots,\square_n)$ with successive centres at
distance at most $M$. Its term is
\[
 S_\omega=(h_0S_0h_0)(C_1S_1h_1)\cdots(C_nS_nh_n),
\]
where the subscript $i$ denotes the cube $\square_i$.
Other sequences give zero because of the support of the commutator.
Its connection variables lie in
$\bigcup_{i=0}^nW_{\square_i}$.

\begin{proposition}\label{prop:walk-propagator-bounds}
For $p\ge1$ with $q(p)=54\pi C_Va'p/M\le1/2$,
\begin{align}
 \sup_z\sum_{z'}\exp(\gamma d(z,z'))
       \sum_{|\omega|=n}p^{n+1}
           \|\mathbf1_{\Delta_z}S_\omega\mathbf1_{\Delta_{z'}}\|
       &\le16a'pw_k(y)q(p)^n,\label{eq:walk-row-bound}\\
 \sup_z\sum_{z'}\exp(\gamma d(z,z'))
       \sum_{|\omega|=n}p^{n+1}
           \|\mathbf1_{\Delta_z}\delta_{\alpha,A}S_\omega
                            \mathbf1_{\Delta_{z'}}\|
       &\le152a'pq(p)^n.\label{eq:walk-holder-bound}
\end{align}
The norms are sup operator norms. At $k=0$ the first bound holds
also with the supremum of the normalized kernel over the two clusters.
In particular, $S_y(A)$ is analytic on $\widetilde{\mathcal G}_k$ and
\begin{equation}\label{eq:global-propagator-bounds}
 N_\gamma(S_y(A))\le32a'w_k(y),\qquad
 N_\gamma(\delta_{\alpha,A}S_y(A))\le304a'.
\end{equation}
These constants are independent of the level, the side, the running
mass in $[0,1]$, and the connection in the stated domain.
\end{proposition}

\begin{proof}
The first factor in a walk has norm at most $2a'w_k(y)$.
Every additional factor has norm at most $2\pi C_Va'/M$ by
\eqref{eq:local-inverse-bounds} and \eqref{eq:cutoff-commutator-bound}.
For a fixed output block, at most eight initial cutoffs contribute;
there are at most $27$ choices at each further step. Apply
submultiplicativity row by row and sum these $8\cdot27^n$ bounds.
This proves \eqref{eq:walk-row-bound}.

For the covariant quotient, only the factor at the left end of the
product is differentiated. The product identity is
\[
 \delta_{\alpha,A}(h_\square v)(x,x')
 =h_\square(c)\delta_{\alpha,A}v(x,x')
   +d_X^{-\alpha}(h_\square(c')-h_\square(c))U_A(x,x')v(x').
\]
The first contribution has localized norm at most $18a'$.
The second has coefficient at most $3\pi\exp(1)/M$, and evaluation
at the second point costs at most $27\exp(\gamma)$. With the stated
choice of $M$ this contribution is at most $a'$. Thus
$N_\gamma(\delta_{\alpha,A}h_\square S_\square)\le19a'$.
For intracluster pairs the second term vanishes, since $c=c'$.
There are again at most eight possible initial cubes, including when
the second point of a short pair enters the cutoff support. The count
$8\cdot27^n$ now gives \eqref{eq:walk-holder-bound}.

At level zero keep the kernel norm in the rightmost local inverse;
the other factors are controlled in operator norm. The same
submultiplicative row estimate proves the kernel assertion.
Summing the geometric bounds with $p=1$ gives
\eqref{eq:global-propagator-bounds}. For a fixed torus the sum of the
operator norms over all walks is finite; its tail of lengths at least
$j$ is bounded by $4a'(\ell/M)^3 2^{-j}$. Each walk term is analytic
in the finitely many connection variables it reads. The uniform tail
bound therefore proves analyticity of the sum. The local bounds in
\eqref{eq:walk-row-bound} contain no factor depending on $\ell$.
\end{proof}


\begin{proof}[Proof of \Cref{thm:chamber-locality}]
Choose $\gamma,a'$ as in \eqref{eq:zero-row-bounds}.
The constants in the free estimates depend only on $\rho,L,\alpha$.
The two charge conditions give \eqref{eq:local-inverse-bounds},
and $q(1)\le1/2$ gives the convergent parametrix.
\Cref{prop:walk-propagator-bounds} supplies the analytic inverse and
both asserted row estimates, including the normalized kernel estimate
at level zero.
\end{proof}
